\chapter{收益从哪儿来：三种钱和一个残酷等式}
\chaptermeta{20}

\begin{ledetext}
你账户里多出来的那几万块，来路一共只有三条。分不清是哪一条，你就永远不知道自己是靠本事还是靠运气。
\end{ledetext}

\section{钱只有三个来路}

上一章那张表做完了，接下来的问题很朴素：怎么让净值那一栏变大。

抛开所有花哨的说法，你从市场上拿走的每一块钱，只可能来自三个地方。

\begin{flowlist}
  \flowitem{01}{经济增长赚的钱}{企业卖出东西、赚到利润、拿一部分分红、剩下的再投进去扩产。这块钱是新造出来的，蛋糕变大了，所有股东按份额一起分。它是正和的。来源说到底只有两样：生产率和人口。}
  \flowitem{02}{别人口袋里的钱}{你 8 块买 12 块卖，赚了 4 块。这 4 块不是凭空长出来的，它对应着某个人 12 块买了、8 块割了。蛋糕没变大，只是重新切了一遍。它是零和的。}
  \flowitem{03}{被摩擦吃掉的钱}{管理费、托管费、申赎费、佣金、买卖价差、印花税、大单进出时把价格打歪造成的冲击成本。这块钱既不来自增长也不来自对手方，它直接离开了所有参与者，去了基金公司、券商和国库。它是负和的，而且每年都收，不管你赚不赚。}
\end{flowlist}

三条路，性质完全不同。第一种可以持续，第二种不能保证，第三种保证发生。

\begin{pullquote}{}
投资这件事，最有用的一个问题是：我赚的这笔钱，到底是哪一种？
\end{pullquote}

\begin{egbox}{把一年的收益拆开看}
老赵在合肥开了家做汽车线束的小厂，业余炒股。2025 年他证券账户赚了 \textbf{4.27 万}，逢人就说自己波段做得好。

拆开：他常年拿着的 22 万宽基指数基金，那年净值涨加分红，贡献 \textbf{2.01 万}（第一种钱）。他另外用 15 万本金做了 11 次短线，来回赚了 \textbf{2.89 万}（第二种钱，对面是那 11 组卖给他和买他手上货的人）。全年佣金、印花税、基金申赎费合计 \textbf{6,300 元}（第三种钱，走了）。

2.01 + 2.89 − 0.63 = 4.27 万。

问题来了：那 2.89 万里，有多少是本事，有多少是运气？11 次样本，答不出来。同样的操作再来五年，如果还是稳定赚，才有一点点说服力。\textbf{第一种钱可以复制，第二种钱只能事后统计。}老赵把两者混在一起，得出了"我会做波段"这个结论。

\end{egbox}

\section{一个不需要任何假设的等式}

1991 年，威廉·夏普写了一篇一页半的东西，叫《主动管理的算术》。全文没有一个统计检验，没有一个假设，只有加法和除法。它的结论至今没人能反驳。

推理是这样走的。

\begin{chainflow}
\chainnode{所有人的持仓加起来就是整个市场}\chainarrow \chainnode{被动投资者按市值权重持有}\chainarrow \chainnode{剩下的主动投资者合起来也必然是市值权重}\chainarrow \chainnode{所以扣费前，主动的平均收益 = 市场收益}\chainarrow \chainnode{主动的成本更高}\chainarrow \chainnode{扣费后，主动的平均收益必然低于指数}
\end{chainflow}

关键在第三步。市场上的股票总共就那么多股，被动那一堆按市值比例拿走了一块，剩下的所有股票落在主动投资者手里，这一堆的构成按定义也是市值权重的。它们不可能整体上"超配"某只股票，因为没有第三方能提供额外的股份。

所以在扣掉成本之前，全体主动投资者作为一个整体，收益率和市场一模一样。不是接近，是相等。然后他们比被动投资者多付了研究费、管理费、换手成本。

减一个正数，只能变小。

\begin{egbox}{一个只有三个人的市场}
某个村子里只有一家公司，总股本 1000 股，年初每股 10 元，市值 1 万元。全村只有老张、老王、老李三个股东。

年底每股 11 元。整个市场涨了 10\%，这是公司经营的结果，跟谁买谁卖没关系。

这一年他们互相买卖了很多次。老张在 11.5 元卖给老王 200 股，赚了个差价；老王吃了亏。但年底一算：三个人合计还是持有 1000 股，合计价值 1.1 万元，\textbf{三人合起来的收益率就是 10\%，一分不多一分不少}。老张多赚的每一分，都是老王少赚的那一分。

再算摩擦。这一年他们交易了 8 次，卖出方的印花税千分之一，双边佣金每笔最低 5 元，合计 \textbf{217 元}离开了村子，进了券商和国库。

三人的真实合计收益是 1000 − 217 = \textbf{783 元}，年化 7.83\%。而村里那个从头到尾没动过、也没人给他讲行情的老李，如果一直满仓不动，他拿到的就是 10\%。

\end{egbox}

\begin{deepbox}{进阶：可以跳过}
这个结论最狠的地方在于它\textbf{不依赖任何关于市场是否有效的假设}。哪怕市场荒谬到人人都在拿彩票号码选股，这个算术照样成立，因为它只用到了"所有人的持仓加起来等于市场"这一条，而这一条是定义，不是观察。

你可以说"这次不一样"，但你得先反驳加法。

它同时也说明了另一件事：它讲的是平均数，不是每个人。一定有人长期跑赢，只是他赢的每一块钱，都必须来自另一个主动投资者输掉的那一块。零和游戏里赢家的存在，恰恰要求输家的存在。\hlmark{你在坐上牌桌之前，最好先想清楚自己打算当哪一个。}

\end{deepbox}

\section{长期收益的真正来源，是你替别人难受}

如果第二种钱靠不住、第三种钱一直在流出，长期回报到底从哪儿来？

答案是第一种钱，而它有一个更精确的名字：\term{风险溢价}{risk premium}\index{015@风险溢价}。

逻辑很朴素。有人手里有钱，但不想承受本金腰斩、不想把钱锁十年、不想承担对方还不上的可能、不想面对"急用钱时卖不掉"。这些不想，都得有人来替他们承担。替他们承担的人，收一笔报酬。

\begin{tablewrap}
\begin{xltabular}{\linewidth}{>{\raggedright\arraybackslash}p{0.20\linewidth}XX}
\toprule
\thead{溢价} & \thead{你替别人承担了什么} & \thead{2026 年的参照} \\
\midrule
\textbf{股权风险溢价} & 企业可能破产、利润可能下修、股价可能腰斩，而且不知道多久回来 & 历史长期在 3–5 个百分点，具体数字学界吵了三十年，至今没定论 \\
\textbf{期限溢价} & 把钱锁上十年，中间利率一动你就浮亏 & 美国 10 年期国债收益率 2026 年 8 月底升至 2025 年 1 月以来最高 \\
\textbf{信用溢价} & 借钱的那一方可能还不上 & 私募信贷借款利率 8\%–12\%，比银行贷款高 200–400 个基点 \\
\textbf{流动性溢价} & 想卖的时候没有买家，或者只能打折卖 & 私募信贷的二顺位与夹层 12\%–16\%，代价是几年不能赎回 \\
\bottomrule
\end{xltabular}
\end{tablewrap}

\begin{pullquote}{}
你赚的不是"聪明"的钱，是"承担别人不愿意承担的东西"的补偿。
\end{pullquote}

这个视角一旦装上，很多事情当场想通。为什么长期看股票比债券收益高？不是因为股票有魔力，是因为持有股票的过程更难受，2008 年跌 70\% 的时候你得拿得住，而拿不住的人把补偿让给了拿得住的人。为什么私募信贷能给你 10\%？因为它有三年不让你赎回，那多出来的几个点是你为"不能反悔"收的钱。

反过来推也成立：\textbf{一个既不用锁定期限、又没有波动、还能给你远高于无风险利率回报的东西，逻辑上不存在。}它要么在某处藏了风险没告诉你，要么它根本就是第 23 章要拆的那种东西。

\section{阿尔法这个词已经被用坏了}

行业里把收益拆成两块。跟着市场一起动的那部分叫\term{贝塔}{beta}\index{003@贝塔}，市场之外多出来的那部分叫\term{阿尔法}{alpha}\index{001@阿尔法}。前者理论上很便宜就能买到，后者才值得付管理费。

麻烦在于"市场"的定义一直在扩张。

1964 年那版模型里，市场只有一个变量。1993 年 Fama 和 French 加进了规模和价值两个\term{因子}{factor}\index{075@因子}，一大批原本被称作阿尔法的收益当场被重新归类为贝塔：那些基金经理不是选股厉害，他们只是长期偏爱小盘和低估值的股票。后来动量、质量、盈利、投资又陆续被加进来，每加一个，就有一批"能力"被翻译成"暴露"。

我给一个判断：\textbf{2026 年市面上被包装成阿尔法的东西，大部分是还没被识别出来的贝塔。}不是造假，是分不清。分不清的代价，是你按阿尔法的价格买了贝塔。

更麻烦的是因子本身在衰减。

想象一家藏在巷子里的面馆，好吃又便宜，因为没什么人知道。有人发了篇测评，三个月后门口排两小时，老板涨了五块钱。面还是那碗面，性价比没了。

因子就是这么死的：论文发表，做成产品，几千亿资金照着同一套信号买进去，便宜的东西不再便宜。往下推一步，结论有点残酷：\hlmark{任何一个被广泛讨论的策略，等你听说它的时候，它最好的那几年大概率已经过去了。}这不是谁在使坏，这就是价格干活的方式。

拥挤还带来第二个后果。所有人用同一套信号，进场时互相抬价，出场时互相踩踏。因子的超额收益变薄了，波动却变大了。这是 2018 年之后量化行业反复经历的事情。

\section{三十年的过路费}

这一节是全章最有杀伤力的部分，只有一道算术题。

本金 30 万，年化毛收益 7\%，投 30 年不动。唯一的变量是费率。

\begin{tablewrap}
\begin{xltabular}{\linewidth}{>{\raggedright\arraybackslash}p{0.20\linewidth}Xrr}
\toprule
\thead{情形} & \thead{净年化} & \thead{30 年后终值} & \thead{被费用拿走} \\
\midrule
\textbf{不收任何费用（理论值）} & 7.00\% & 228.4 万 & 0 \\
\textbf{宽基指数基金，费率 0.15\%} & 6.85\% & 218.9 万 & 9.5 万 \\
\textbf{主动型基金，费率 1.5\%} & 5.50\% & 149.5 万 & 78.9 万 \\
\bottomrule
\end{xltabular}
\end{tablewrap}

\begin{barfig}
  \barrow{零费用}{100}{228.4 万}
  \barrow{费率 0.15\%}{96}{218.9 万}
  \barrow{费率 1.5\%}{65}{149.5 万}
\end{barfig}

218.9 万和 149.5 万，差 \textbf{69.4 万}，占了高费率那一栏终值的将近一半，占了低费率终值的三成二。

把 69.4 万换个说法：它是本金 30 万的 2.3 倍。你投进去 30 万，管理人从中拿走了 69 万。

合同上写的差距只有 1.35 个百分点。1.35\% 看起来像个零头，很多人在签字的时候连问都不问。

\begin{pullquote}{}
这不是收益率的差别，这是复利被人在门口收了三十年过路费。
\end{pullquote}

要说公道话：如果一只基金能持续跑赢指数 3 个百分点，1.5\% 的费率是便宜的。问题是持续跑赢 3 个百分点的人极少，而且你在买入的那一刻无法辨认他是谁。\hlwarm{这段不构成投资建议}，它只是在说，费率是你事前唯一能确定知道的数字，收益率不是。

\section{为什么亏钱的总是散户}

下面几条，我自己全中过。写出来不是为了嘲笑谁。

\textbf{赚了就跑，亏了就扛。}这个毛病学名叫\term{处置效应}{disposition effect}\index{009@处置效应}。涨 8\% 的时候想落袋为安，跌 30\% 的时候说"反正没卖就不算亏"。后半句在会计上是错的：你的账户余额已经变了，没卖唯一改变的是你的心情。结果是你手里长期留着最差的那些，卖掉的全是最好的那些。

\textbf{过度自信。}各国的调查里，大约八成司机认为自己车技高于平均水平。这在数学上做不到。股票账户里的版本是：亏了怪市场，赚了归自己。

\textbf{买你看到的，不是买你研究的。}热搜、涨停板、朋友圈里的收益截图。这里有一个不对称：没人晒亏损截图。所以你看到的样本从第一眼开始就是歪的，你以为自己在观察世界，其实在观察幸存者。

\textbf{手一直在动。}换手率越高，成本越高，收益越低。这一条在各国的散户账户研究里，几乎是唯一一个稳定复现的结论。它跟你选股水平无关，它是算术。

这四条有个共同点：它们全都发生在你觉得自己很理性的时候。

\section{2026 年，你站在哪儿}

\begin{databox}{2026 数据}
\textbf{全球被动投资规模已首次超越主动管理。}这件事花了五十年。夏普那篇一页半的文章写于 1991 年，市场用了三十五年才开始按它说的做。

\textbf{中国公募基金规模约 38 万亿元}（2026 年初数据）。这个盘子已经大到，公募之间互相博弈的部分远大于它们从散户手里赚走的部分。

\textbf{A 股：2026 年 8 月 31 日收盘，上证指数 3986.30 点，当月上涨 4.02\%，两市成交约 2.13 万亿元。}市场的核心逻辑正从"估值修复"转向"盈利驱动"。

\end{databox}

最后那句话如果翻译成人话，正好落回本章开头。

\begin{talkbox}
  \talkline{研报}{"市场逻辑正从估值修复转向盈利驱动。"}
  \talkline{翻译}{过去这一段，股价涨是因为大家愿意为同样的一块钱利润多付几倍钱（第二种钱，重新定价，有人接手就涨，没人接手就停）。往后要涨，得靠那一块钱利润本身变成一块二（第一种钱，蛋糕真的变大）。}
  \talkline{所以}{估值能修复到哪儿，没人知道，那取决于别人愿意付多少。盈利能不能兑现，财报里会有答案，虽然要等。}
\end{talkbox}

\begin{mythbox}{常见误解}
\mvfrow{误解}{"投资就是找到会涨的股票。"}
\mvfrow{事实}{拉长到十年以上看，宽基指数的累计收益里绝大部分来自企业盈利增长和分红再投入，估值变化的贡献要小得多。"找到会涨的股票"这句话里藏着一个没说出口的前提：会有人在更高的价位从你手里接走。\textbf{研究"哪家公司能赚钱"和研究"下个月谁愿意出更高价"，是两门完全不同的手艺。}很多人以为自己在做前者，账户的换手率出卖了他们。}
\end{mythbox}

\begin{mythbox}{常见误解}
\mvfrow{误解}{"基金经理是专业的，肯定比我强。"}
\mvfrow{事实}{他在信息处理和纪律上多半确实比你强。但"比你强"和"能持续跑赢同类平均、并且把 1.5\% 的费用赚回来"，是两个命题。第一个成立不代表第二个成立。\\ 真正卡住你的是第三件事：\textbf{即使确实存在少数能持续跑赢的人，你在事前挑不出他。}用过去三年排名去挑，前四分之一的基金在下一个三年继续留在前四分之一的比例，长期统计下来接近抛硬币。而销售页面上给你看的，永远是过去三年的排名。}
\end{mythbox}

\begin{warnbox}{小心}
本章讲的全是机制，不是标的。哪个指数会涨、哪只基金明年好、A 股 3986 点是高是低，这些问题本书一个都不回答，也建议你对任何回答得很干脆的人保持警惕。\hlwarm{对机制可以下判断，对价格不能。}

\end{warnbox}

\begin{takeaway}{本章一句话}
你赚的每一块钱只有三个来路：经济增长、别人的口袋、以及正在被摩擦吃掉。第一种可以持续，第二种是零和的，第三种保证发生。

\begin{itemize}
  \item 扣费前主动投资者的平均收益必然等于市场，扣费后必然低于指数。这是恒等式，不是观点。
  \item 长期收益来自承担别人不愿承担的东西：波动、期限、违约、卖不掉。
  \item 30 万投 30 年，费率 0.15\% 和 1.5\% 的终值差 69.4 万，是本金的 2.3 倍。
  \item 费率是你事前能确定的唯一数字。
\end{itemize}

\end{takeaway}

\begin{bridgebox}
\textbf{下一章}既然长期收益主要来自承担风险，那有没有办法在不降低收益的前提下少担一点？有，而且只有这一个。
\end{bridgebox}

